\documentclass[11pt,oneside]{book}
\usepackage[
    paperwidth=6in,
    paperheight=9in,
    top=0.75in,
    bottom=0.75in,
    left=0.6in,
    right=0.6in,
    twoside=false
]{geometry}
\input{../pri-end/T0_preamble_En.tex}
% T0_preamble_patches.tex — LOKALER PLATZHALTER
% Im Repo durch die offizielle Version ersetzen.

\usepackage{graphicx}

\makeatletter
\renewcommand{\thechapter}{\arabic{chapter}}
\makeatother
\raggedbottom
\allowdisplaybreaks
\hfuzz=5pt
\vfuzz=5pt
\emergencystretch=2em

\setcounter{tocdepth}{0}

\begin{document}
\frontmatter
\pagestyle{fancy}

% ── Title page ──────────────────────────────────────────────
\begin{titlepage}
  \centering\vspace*{1.5cm}
  {\Huge\bfseries\color{blue!70!black}
    Dark Matter — an Illusion}\\[0.6cm]
  {\Large\bfseries\color{blue!40!black}
    Why Galaxies Work Without Invisible Mass}\\[1cm]
  {\LARGE\itshape Fundamental Fractal Geometric\\[0.2cm]
    Field Theory (FFGFT) — Galactic Physics}\\[0.5cm]
  {\large\itshape A Popular Science Overview}\\[2cm]
  {\large Johann Pascher}\\[0.3cm]
  {\normalsize ORCID: 0009-0000-6518-4064}\\[0.3cm]
  {\normalsize\url{https://github.com/jpascher/T0-Time-Mass-Duality}}\\[1.5cm]
  {\normalsize September 2026}
  \vfill{\small Version 1.0 --- Popular Science Edition}
\end{titlepage}

\newpage\thispagestyle{empty}%
{\small\noindent\textbf{Keywords:} Dark matter $\cdot$
rotation curves $\cdot$ galaxies $\cdot$ MOND $\cdot$
gravitational lensing $\cdot$ Bullet Cluster $\cdot$ FFGFT $\cdot$
inertia regime\\[1em]
\noindent Johann Pascher: \textit{Dark Matter — an Illusion.
Why Galaxies Work Without Invisible Mass. FFGFT.}\\
First edition, September 2026.\\[0.6em]
\noindent Licence: Creative Commons BY 4.0}

\tableofcontents
\mainmatter

% ════════════════════════════════════════════════════════════
\chapter{The Mystery of Flat Curves}
% ════════════════════════════════════════════════════════════

It was a simple measurement with an uncomfortable answer.

In the 1970s, Vera Rubin observed something that physics has
not let go of since: the stars at the edge of a galaxy move
too fast. Far too fast. According to the law of gravity —
the same law that keeps planets orbiting the Sun — stars
far from the centre should move more slowly than those
closer in. Just as Neptune orbits the Sun more slowly than
Earth does.

But they do not.

The rotation curves of galaxies — the curve showing how fast
stars at different distances from the centre orbit — flatten
out at the edges. The velocity stays constant as far as one
can measure. It is as if all planets in our solar system
orbited the Sun at the same speed, no matter how far away
they are.

The standard explanation is: there must be more mass than we
see. An invisible, lightless mass — dark matter — must
envelop the galaxies and strengthen the gravitational pull.
Without it, the stars at the edge would fly away.

This explanation sounds reasonable. But it comes at a price:
dark matter has never been directly detected. Not in particle
accelerators, not in underground detectors, not through any
other experiment. We postulate 26\,\% of the entire universe
to consist of a substance we do not know, cannot see, and
cannot grasp.

This book asks a different question: What if galaxies are
not misbehaving — but our physics simply stops working at
their edges?

FFGFT — Fundamental Fractal Geometric Field Theory — gives
a precise answer. It needs no new particles. It needs no
invisible mass. It needs only a single number:
$\xi \approx \tfrac{4}{3} \times 10^{-4}$.

% ════════════════════════════════════════════════════════════
\chapter{What Newton Really Says — and Where He Stops}
% ════════════════════════════════════════════════════════════

Newton's law of gravity is one of the most reliable formulas
in physics. On Earth, in the solar system, in computing
satellite trajectories — it works with extraordinary
precision.

But it carries a hidden assumption: that mass and inertia
always bear the same relationship to each other under all
circumstances. Inertia is what prevents a body from changing
its velocity — the greater the inertia, the more force is
needed to accelerate it.

Newton wrote: $F = m \cdot a$. Force equals mass times
acceleration. This seems self-evident. But what if this
relationship only holds for large accelerations?

In everyday life, accelerations are large. A falling object
on Earth experiences $g \approx 10~\text{m/s}^2$. A
satellite in Earth orbit perhaps a thousandth of that. But
a star at the outermost edge of a galaxy experiences an
acceleration of only $a \sim 10^{-10}~\text{m/s}^2$ —
one ten-billionth of Earth's surface gravity.

Have we ever tested whether Newton's law holds in this regime?

The honest answer is: no. In the laboratory, in the solar
system, on Earth — we have confirmed it thousands of times.
In the galactic deep-field regime, we have never tested it
directly.

FFGFT says: different physics applies there. Not because a
new force appears, but because inertia itself becomes
nonlinear.

% ════════════════════════════════════════════════════════════
\chapter{The Fundamental Relation: Time and Mass as Reciprocals}
% ════════════════════════════════════════════════════════════

The core of FFGFT is a single formula:

\[
  \tilde{T} \cdot m = 1 .
\]

Intrinsic time and mass are reciprocals of each other.
Where mass is large, time passes slowly. Where mass is
small, time passes quickly. This is not a metaphor — it is
a geometric property of the field.

From this relation it follows: a body with very small mass
inhabits a time regime that differs fundamentally from that
of a heavy body. And since inertia is connected to mass and
time, inertia itself changes.

This sounds abstract. But it has a very concrete consequence
for galaxies.

At the edge of a galaxy, the mass of visible matter is low,
and the acceleration is very small. Both effects appear
together: small mass and small acceleration. FFGFT says that
in this regime, Newton's formula $F = ma$ no longer holds
exactly. Inertia becomes smaller than Newton predicts —
stars respond to gravity more sensitively than expected.

This is the effect we see in rotation curves — without dark
matter.

% ════════════════════════════════════════════════════════════
\chapter{The Transition Scale $a_0$ — a Boundary from First Principles}
% ════════════════════════════════════════════════════════════

For a theory to be scientific, it must predict precisely
where an effect begins. FFGFT does this.

There is an acceleration scale $a_0$ below which the
Newtonian regime ends and the galactic regime begins.
This scale is not a free number to be fitted — it follows
directly from the single parameter $\xi$ of the theory:

\[
  a_0 = \frac{c^2 \xi^{10}}{4\,\bar{\lambda}_e}
    \approx 1{.}03 \times 10^{-10}~\text{m/s}^2 .
\]

Here $\bar{\lambda}_e = \hbar/(m_e c)$ is the reduced
Compton wavelength of the electron — a known, measurable
quantity. The factor $\xi^{10}$ carries the entire scale
structure: ten powers of the fundamental constant connect
the quantum scale of the electron to the scale of galactic
acceleration.

This is remarkable. This value of $a_0$ had been known
empirically in astrophysics for decades — Mordehai Milgrom
determined it from observations in 1983, without being able
to derive it. FFGFT derives it from first principles, with
no free parameter.

Below $a_0$, the modified force law applies:

\[
  a = \sqrt{a_N^2 + a_N\,a_0} ,
\]

where $a_N$ is the Newtonian acceleration from the visible
mass. For large accelerations ($a_N \gg a_0$) this reduces
to Newton: $a \to a_N$. For small accelerations
($a_N \ll a_0$):

\[
  a \to \sqrt{a_N\,a_0} \qquad\Longrightarrow\qquad
  v^4 = G\,M_b\,a_0 .
\]

This is the \emph{baryonic Tully-Fisher relation} — one of
the most robust empirical rules in galactic physics,
observed for forty years, yet never derived from first
principles by the standard model. FFGFT derives it.

% ════════════════════════════════════════════════════════════
\chapter{The Test: DDO 154 and the Milky Way}
% ════════════════════════════════════════════════════════════

A theory must be measured against observations. Three
galaxies of different types are particularly instructive.

\textbf{DDO 154 — the cleanest test.}
This dwarf galaxy is gas-dominated: almost all of its
visible mass consists of hydrogen gas, whose quantity is
known very precisely from radio astronomy. There is little
uncertainty about the visible mass — this makes DDO 154
the ideal test case.

The Newtonian prediction for the outer rotation velocity
is about 18.5~km/s. The observation shows 47.0~km/s —
more than twice as large.

The FFGFT prediction from the modified force law: 47.1~km/s.

The ratio of prediction to observation: 1.001 — agreement
to one part per thousand.

\textbf{NGC 3198 — a spiral galaxy.}
Here the baryonic mass is less precisely known (uncertainty
$\sim 30\,\%$). The Newtonian expectation: 50.4~km/s,
observed: 150~km/s. FFGFT predicts 121~km/s — a deviation
of 18\,\%, fully within the mass uncertainty.

\textbf{The Milky Way.}
Newton expects 122~km/s in the outer disc, measurements
give around 220~km/s, FFGFT predicts 180~km/s. Again,
the deviation lies within the stellar mass uncertainty.

\medskip
Dark matter explains these curves too — but only by fitting
a separate halo profile for each galaxy. Those are free
parameters, one per galaxy. FFGFT has not a single free
parameter: $a_0$ follows from $\xi$, which itself carries
all of physics.

% ════════════════════════════════════════════════════════════
\chapter{The Bullet Cluster — the Seemingly Strongest Argument}
% ════════════════════════════════════════════════════════════

Anyone who argues against dark matter eventually encounters
the same objection: the Bullet Cluster.

In 2006, Douglas Clowe and colleagues published an
observation considered decisive evidence for dark matter.
Two galaxy clusters had collided. The hot gas — by far the
largest visible mass component — was slowed by the collision
and left behind. The galaxies themselves passed through.

From gravitational lensing measurements, the true location
of the mass was determined: not at the gas, but at the
galaxies. The gas was slowed; the mass was not. Therefore,
a collision-free invisible mass must be present — dark matter.

This is a compelling argument. But FFGFT answers it without
dark matter.

The modified force law acts on all mass — visible and
interacting alike. What matters is the mass ratio between
galaxies and gas in a cluster. The galaxies carry roughly
ten times more mass than the gas. During the collision, the
gas is slowed by radiation pressure and electromagnetic
interaction — the galaxies are not.

In a finite-element simulation of the collision geometry
(using the parameters of Clowe et al.\ 2006), FFGFT gives:
the gravitational lensing peak lies 10~kpc from the
galaxies and 240~kpc from the gas centre. This is exactly
what is observed.

The Bullet Cluster is not evidence for dark matter. It is
evidence that mass and gas respond differently in a
collision. This has always been known — and in FFGFT it
follows directly from the geometry.

% ════════════════════════════════════════════════════════════
\chapter{Dark Energy — a Construct of the Wrong Interpretation}
% ════════════════════════════════════════════════════════════

Alongside dark matter, the standard model knows a second
unknown: dark energy. It is supposed to drive the
accelerated expansion of the universe and is said to
account for 69\,\% of the entire universe.

Together, the two unknown components amount to 95\,\% of
the universe. Only 5\,\% is ordinary matter that we know
and understand.

FFGFT needs neither dark matter nor dark energy — but for
different reasons.

Dark matter, as shown, is not a substance: it is a regime
in which inertia becomes nonlinear. The effects are real —
rotation curves are real, gravitational lensing is real —
but they require no new ingredient.

Dark energy disappears for a different reason: there is no
expansion to accelerate. In FFGFT the universe is static.
The redshift of distant galaxies arises not because space
is stretching, but because light loses energy along its
path through the $\xi$-field:

\[
  1 + z = e^{\xi\,x} .
\]

For small distances this is indistinguishable from a Doppler
effect. For large distances it makes the same predictions
as the standard model — because the data are degenerate:
they cannot in principle distinguish between expansion and
geometric path lengthening.

The cosmological constant $\Lambda$ — the mathematical
shorthand for dark energy — has no referent in a static
universe. It is a bookkeeping quantity of the expansion
interpretation, not a physical object.

% ════════════════════════════════════════════════════════════
\chapter{What the Data Can Actually Decide}
% ════════════════════════════════════════════════════════════

A fair question: if FFGFT can explain the same observations
as the standard model — why switch?

The answer lies in the structure of the theories, not in
individual data points.

The standard cosmological model ($\Lambda$CDM) describes
cosmology with six free parameters, all empirically
determined: Hubble constant, baryon density, dark matter
density, spectral index, fluctuation amplitude, optical
depth. The particle physics standard model adds 19 more.
In total, over 25 numbers, none derived from first
principles.

FFGFT has a single parameter: $\xi = \tfrac{4}{3}
\times 10^{-4}$. From it follow:
\begin{itemize}
  \item the gravitational constant $G$,
  \item the Hubble constant $H_0 = 66{.}82~\text{km/s/Mpc}$,
  \item the acceleration scale $a_0 = 1{.}03 \times 10^{-10}~\text{m/s}^2$,
  \item the CMB temperature $T_\text{CMB} = 2{.}725~\text{K}$,
  \item the baryonic Tully-Fisher relation,
  \item the lepton mass hierarchy.
\end{itemize}

The same $\xi$ anchors five different sectors of physics —
particle physics, cosmology, galactic dynamics, gravitational
lensing, quantum mechanics. That is not calibration: it is
cross-anchoring.

The cosmological constant $\Lambda$, by contrast, works
nowhere outside its own fit. That is the difference between
a theory and a bookkeeping entry.

\medskip

Where can one decide today? The classical pillars —
supernovae, baryonic acoustic oscillations, CMB spectrum —
are degenerate: both models produce the same mathematical
curves. But there are observations that discriminate:

\begin{itemize}
  \item \textbf{Wavelength dependence of redshift:}
    FFGFT predicts strictly achromatic redshift.
    Any demonstrated wavelength dependence would falsify
    the model.
  \item \textbf{Direct dark matter detection:}
    Any direct detection of a dark matter particle —
    in accelerators, detectors, or cosmic sources —
    would show that a new substance genuinely exists.
  \item \textbf{Euclid and DESI:}
    The new galaxy surveys will measure the growth of
    cosmic structure with unprecedented precision.
    FFGFT makes concrete predictions for the power
    spectrum without dark matter halo corrections.
\end{itemize}

% ════════════════════════════════════════════════════════════
\chapter{An Honest Assessment}
% ════════════════════════════════════════════════════════════

FFGFT explains galaxy rotation curves without dark matter.
It derives the acceleration scale $a_0$ from first
principles. It passes the Bullet Cluster test. It explains
why dark energy is not a physical object.

This is not a complete theory of gravity. Open questions
remain:

\textbf{The K3 forward proof} — the mathematical derivation
of the geometry from which the modified force law follows
is structurally motivated but not yet fully proved.

\textbf{14\,\% deviation at $a_0$:} The theoretically
derived value lies about 14\,\% above the midpoint of the
empirical range. However, the literature values for $a_0$
scatter between $1{.}1$ and $1{.}3 \times 10^{-10}~\text{m/s}^2$
— the FFGFT value lies within this range.

\textbf{Cluster physics:} For large galaxy clusters with
significant mass uncertainty, the predictions are less
precise than for gas-dominated dwarf galaxies.

\medskip

What is clear: assuming that 95\,\% of the universe
consists of unknown substances that have never been directly
detected is not a smaller assumption than assuming that
inertia becomes nonlinear at very small accelerations.

The difference: the former assumption is not falsifiable —
one can always add more dark matter when an observation
fails to fit. The latter has a precise, distinguished scale
$a_0$ that follows from the structure of the theory and
manifests in observations.

Dark matter is not wrong. It might exist. But it is not
necessary. And a physics that manages without it deserves
serious scrutiny.

% ════════════════════════════════════════════════════════════
\chapter{What Comes Next}
% ════════════════════════════════════════════════════════════

The next generation of observations will decide.

The ESA Euclid telescope is surveying billions of galaxies
out to redshifts of $z \sim 2$. The James Webb Space
Telescope is showing us galaxies in an early universe that,
according to the standard model, should not yet exist —
but fit comfortably in the static FFGFT universe. DESI is
measuring baryonic acoustic oscillations with record
precision.

For FFGFT, concrete predictions remain open:

Achromatic redshift — no wavelength dependence across the
entire observable range — is a core test. Current
precision observations already set tight limits; FFGFT is
so far compatible.

The mass distribution in galaxy clusters without dark
matter halos: FFGFT cluster profiles differ subtly from
NFW profiles of standard dark matter models. The resolution
of future lensing surveys could distinguish them.

The baryonic Tully-Fisher relation holds in FFGFT exactly
— not as a statistical trend with scatter, but as a
geometric consequence. Deviations would be a direct
indication of dark matter as a genuine substance.

\medskip

Whether dark matter exists or is an illusion of the wrong
physics is one of the greatest open questions in science.
It is decidable. The theory that answers it will not be
found through belief — but through predictions that either
prove themselves or fail.

FFGFT stands ready to be tested.

\backmatter

% ── Sources ──────────────────────────────────────────────────
\chapter*{Sources and Foundation Documents}
\addcontentsline{toc}{chapter}{Sources}

This book is based on the FFGFT corpus in the GitHub
repository \url{https://github.com/jpascher/T0-Time-Mass-Duality}.
The following documents are the primary sources:

\begin{itemize}
  \item \textbf{Doc.~308} — The cosmic sector of FFGFT
    (redshift, dark matter, Bullet Cluster, light deflection)
  \item \textbf{Doc.~267} — Cosmological degeneracy:
    $\Lambda$CDM and FFGFT as two readings of the same observations
  \item \textbf{Doc.~309} — The scale-anchor problem
    ($H_0$ derivation, SI anchor, asymmetry of the models)
  \item \textbf{Doc.~025} — T0 cosmology
    (CMB, $\xi$-field effects, static universe)
  \item \textbf{Doc.~004} — T0 model overview
    (fundamental relations, field geometries)
  \item \textbf{Doc.~012} — Gravitational constant from FFGFT
\end{itemize}

All calculations in Chapters~4 and~5 were verified with the
verification scripts
\texttt{ffgft\_308\_p44\_stufeA\_a0\_kette.py} and
\texttt{ffgft\_308\_p44\_stufeC\_bullet\_cluster.py}.

\end{document}
